\documentclass[twocolumn]{aastex631}
\begin{document}
\title{The reionization ionizing-photon-budget shortfall for star-forming galaxies is not robust to systematics at z\~{}6 (o32 systematic corner)}
\author{NebulaMind Lab (autonomous pipeline)}
\affiliation{NebulaMind Lab, \url{https://lab.nebulamind.net}}
\begin{abstract}
Reionization ionizing-photon-budget reconciliation at z\~{}6: star-forming galaxies require f\_esc=0.072 (+0.072/-0.037) to close the budget (Madau-Dickinson SFRD, log xi\_ion=25.5+/-0.15, clumping C=2-5, JWST-SFRD tail), versus indirect-proxy-inferred f\_esc=0.080 (+0.146/-0.051) from LzLCS O32/beta calibrations. Median delta(required-inferred)=-0.007 dex-frac (16-84\%: -0.147 to +0.073); 46\% of the systematic MC shows a shortfall. Conclusion: the budget CLOSES within the systematic, and the sign flips sign between the O32 and beta calibrations (calibration-driven, not robust). The result is bounded by the xi\_ion x clumping x proxy-calibration systematic, not by statistics. Generated autonomously from public data (jwst) via the NebulaMind Lab runner: a bounded, reproducible, descriptive study.
\end{abstract}
\section{Introduction}
Introduction:
Recent studies have highlighted a potential crisis in reconciling the ionizing photon budget during reionization [Muñoz2024]. This issue arises from discrepancies between the estimated number of ionizing photons produced by star-forming galaxies and the actual requirements for reionizing the universe. To address this, we revisit the reionization-photon-budget using a literature-anchored approach, relying on established values from previous works [Madau2017] to inform our analysis.
\section{Data and method}
Data and Method:
Our calculation is based on the cosmic star formation rate density (SFRD) provided by Madau \& Dickinson's analytic fitting function [Madau2014]. We adopt published values for xi\_ion and O32/beta f\_esc proxy calibrations from LzLCS studies [Chisholm+22, Flury+22; Simmonds+24]. Notably, we do not utilize survey catalog data or observations from JWST, SDSS, or TNG in this analysis. Instead, our method focuses on reconciling the ionizing-photon-budget through a systematic examination of published literature values.
\section{Result}
Result:
Our reconciliation of the reionization ionizing-photon-budget at z\~{}6 reveals that star-forming galaxies require an escape fraction f\_esc=0.072 (+0.072/-0.037) to close the budget, assuming the Madau-Dickinson SFRD, log xi\_ion=25.5±0.15, clumping factor C=2-5, and JWST-SFRD tail. This value is compared to indirect-proxy-inferred f\_esc=0.080 (+0.146/-0.051) derived from LzLCS O32/beta calibrations. The median delta between required and inferred values is -0.007 dex-frac (16-84\%: -0.147 to +0.073), with 46\% of systematic Monte Carlo simulations showing a shortfall.
\begin{figure}\centering\includegraphics[width=\columnwidth]{result.png}\caption{The reionization ionizing-photon-budget shortfall for star-forming galaxies is not robust to systematics at z\~{}6 (o32 systematic corner)}\end{figure}
\section{Caveats}
Caveats:
It is essential to acknowledge the limitations of our approach, which relies on automated selection and uncalibrated measurements from published literature. The accuracy of our result is contingent upon the assumptions and calibrations used in previous studies, such as the Madau-Dickinson SFRD and LzLCS proxy calibrations. Furthermore, our analysis does not account for potential systematic errors or uncertainties inherent in these sources, which may impact the validity of our findings. Additionally, the reliance on a single selection method and lack of direct observational data may introduce biases that are not fully addressed in this study.
\section*{References}
\footnotesize
[Muoz2024] Muñoz et al. (2024). Reionization after JWST: a photon budget crisis?. bibcode 2024MNRAS.535L..37M \par [Davies2021] Davies et al. (2021). The Predicament of Absorption-dominated Reionization: Increased Demands on Ionizing Sources. bibcode 2021ApJ...918L..35D \par [Park2022] Park et al. (2022). Calibrating excursion set reionization models to approximately conserve ionizing photons. bibcode 2022MNRAS.517..192P \par [Duncan2015] Duncan et al. (2015). Powering reionization: assessing the galaxy ionizing photon budget at z < 10. bibcode 2015MNRAS.451.2030D \par [Madau2017] Madau (2017). Cosmic Reionization after Planck and before JWST: An Analytic Approach. bibcode 2017ApJ...851...50M

\end{document}
